\documentclass{article}
\usepackage[T1]{fontenc}
\usepackage{lmodern}
\usepackage{graphicx}
\usepackage{amsmath,amssymb}
\usepackage[a4paper,margin=2cm]{geometry}
\usepackage[absolute,overlay]{textpos}
\setlength{\TPHorizModule}{1mm}
\setlength{\TPVertModule}{1mm}
\usepackage[indonesian]{babel}
\usepackage{hyperref}
\setlength{\emergencystretch}{2em}
% unit: Halaman 46

\begin{document}

% === Halaman 46 (llm) ===

\section*{Counter Mode}

\begin{itemize}
    \item Mode \textit{counter} tidak melakukan perantaian (\textit{chaining}) seperti pada CBC.
    \item \textit{Counter} adalah sebuah nilai berupa blok bit yang ukurannya sama dengan ukuran blok plainteks.
    \item Nilai \textit{counter} harus berbeda dari setiap blok yang dienkripsi. Pada mulanya, untuk enkripsi blok pertama, \textit{counter} diinisialisasi dengan sebuah nilai.
    \item Selanjutnya, untuk enkripsi blok-blok berikutnya \textit{counter} dinaikkan (\textit{increment}) nilainya satu \[ \text{counter} \leftarrow \text{counter} + 1 ). \]

\end{document}
